\documentclass[10pt,a4paper]{amsart}
\usepackage{amsmath,amssymb,amsthm,graphicx}
\usepackage{verbatim}
\newenvironment{reference}{\comment}{\endcomment}
%\newenvironment{reference}{}{}
\newenvironment{slogan}{\comment}{\endcomment}
\newenvironment{history}{\comment}{\endcomment}

% For commutative diagrams we use Xy-pic
\usepackage[all]{xy}

% We use 2cell for 2-commutative diagrams.
\xyoption{2cell}
% Theorem environments.
%
\theoremstyle{plain}
\newtheorem{theorem}[subsection]{Theorem}
\newtheorem{proposition}[subsection]{Proposition}
\newtheorem{lemma}[subsection]{Lemma}

\theoremstyle{definition}
\newtheorem{definition}[subsection]{Definition}
\newtheorem{example}[subsection]{Example}
\newtheorem{exercise}[subsection]{Exercise}
\newtheorem{situation}[subsection]{Situation}

\theoremstyle{remark}
\newtheorem{remark}[subsection]{Remark}
\newtheorem{remarks}[subsection]{Remarks}

\numberwithin{equation}{subsection}

% Macros
%
\def\lim{\mathop{\mathrm{lim}}\nolimits}
\def\colim{\mathop{\mathrm{colim}}\nolimits}
\def\Spec{\mathop{\mathrm{Spec}}}
\def\Hom{\mathop{\mathrm{Hom}}\nolimits}
\def\Ext{\mathop{\mathrm{Ext}}\nolimits}
\def\SheafHom{\mathop{\mathcal{H}\!\mathit{om}}\nolimits}
\def\SheafExt{\mathop{\mathcal{E}\!\mathit{xt}}\nolimits}
\def\Sch{\mathit{Sch}}
\def\Mor{\mathop{\mathrm{Mor}}\nolimits}
\def\Ob{\mathop{\mathrm{Ob}}\nolimits}
\def\Sh{\mathop{\mathit{Sh}}\nolimits}
\def\NL{\mathop{N\!L}\nolimits}
\def\CH{\mathop{\mathrm{CH}}\nolimits}
\def\proetale{{pro\text{-}\acute{e}tale}}
\def\etale{{\acute{e}tale}}
\def\QCoh{\mathit{QCoh}}
\def\Ker{\mathop{\mathrm{Ker}}}
\def\Im{\mathop{\mathrm{Im}}}
\def\Coker{\mathop{\mathrm{Coker}}}
\def\Coim{\mathop{\mathrm{Coim}}}

% Boxtimes
%
\DeclareMathSymbol{\boxtimes}{\mathbin}{AMSa}{"02}

%
% Macros for moduli stacks/spaces
%
\def\QCohstack{\mathcal{QC}\!\mathit{oh}}
\def\Cohstack{\mathcal{C}\!\mathit{oh}}
\def\Spacesstack{\mathcal{S}\!\mathit{paces}}
\def\Quotfunctor{\mathrm{Quot}}
\def\Hilbfunctor{\mathrm{Hilb}}
\def\Curvesstack{\mathcal{C}\!\mathit{urves}}
\def\Polarizedstack{\mathcal{P}\!\mathit{olarized}}
\def\Complexesstack{\mathcal{C}\!\mathit{omplexes}}
% \Pic is the operator that assigns to X its picard group, usage \Pic(X)
% \Picardstack_{X/B} denotes the Picard stack of X over B
% \Picardfunctor_{X/B} denotes the Picard functor of X over B
\def\Pic{\mathop{\mathrm{Pic}}\nolimits}
\def\Picardstack{\mathcal{P}\!\mathit{ic}}
\def\Picardfunctor{\mathrm{Pic}}
\def\Deformationcategory{\mathcal{D}\!\mathit{ef}}
\usepackage{cleveref}
\crefname{originalreference}{original reference}{original references}
\Crefname{originalreference}{Original reference}{Original references}
\makeatletter
\newlabel{situation-shriek}{{0F42}{source}{}{sourcepointer-1}{}}
\newlabel{situation-shriek@cref}{{[originalreference][0][]0F42}{[1][0][]source}{}{}{}}
\makeatother
\crefname{originallemma}{original lemma}{original lemmas}
\Crefname{originallemma}{Original lemma}{Original lemmas}
\makeatletter
\newlabel{lemma-relative-dualizing-fibres}{{0BZZ}{source}{}{sourcepointer-2}{}}
\newlabel{lemma-relative-dualizing-fibres@cref}{{[originallemma][0][]0BZZ}{[1][0][]source}{}{}{}}
\makeatother
\crefname{originallemma}{original lemma}{original lemmas}
\Crefname{originallemma}{Original lemma}{Original lemmas}
\makeatletter
\newlabel{lemma-gorenstein}{{0BFQ}{source}{}{sourcepointer-3}{}}
\newlabel{lemma-gorenstein@cref}{{[originallemma][0][]0BFQ}{[1][0][]source}{}{}{}}
\makeatother
\crefname{originallemma}{original lemma}{original lemmas}
\Crefname{originallemma}{Original lemma}{Original lemmas}
\makeatletter
\newlabel{more-algebra-lemma-lift-bounded-pseudo-coherent-to-perfect}{{0BYP}{source}{}{sourcepointer-4}{}}
\newlabel{more-algebra-lemma-lift-bounded-pseudo-coherent-to-perfect@cref}{{[originallemma][0][]0BYP}{[1][0][]source}{}{}{}}
\makeatother
\title{For a flat morphism Gorensteinness at a point is equivalent to the shriek being invertible nearby}
\author{The Stacks Project Authors}
\date{stacks-project:duality.tex:local-snapshot-20261012; GNU FDL 1.2 or later}
\begin{document}
\maketitle
\noindent\textit{Editorial note.} All delivered mathematical text is the author's own, copied from the pinned Stacks Project chapter file duality.tex; no statement or proof step was rewritten, reordered or removed.\par\medskip
\noindent\textit{Editorial note.} Preamble provenance: the pinned chapter file opens with an input of the shared preamble and carries no preamble text of its own, so the project's single Stacks-wide preamble staged at sources/stacks-duality/preamble.tex is excerpted verbatim under the recipe fixed by LANE\_BRIEF 22.1, as three spans, lines 8-13 (verbatim package and the reference/slogan/history environments), lines 14-18 (Xy-pic setup) and lines 154-217 (the native theorem declarations and mathematical macros). Nothing was added, removed or redefined; the wrapper class and the amsmath/amssymb/amsthm/graphicx base packages are supplied by the assembler.\par\medskip
\noindent\textit{Editorial note.} The counted statement and proof contain no cross reference, no citation and no included file, so the delivered body is self contained.\par\medskip
\noindent\textit{Editorial note.} Canonical location of the counted proposition: Stacks tag 0C08 (sources/stacks-duality/pinned-original-source.tex line 6722).\par\medskip
\noindent\textit{Allowed standard prerequisites (editorial).}\begin{itemize}
\item The shriek of the structure sheaf is a bounded complex with coherent cohomology whose fibre restriction is dualizing. Same chapter, Lemma labelled lemma-relative-dualizing-fibres, cited at proof lines 6723-6727. Makes the derived restriction to the fibre a dualizing complex.
\item The criterion for Gorensteinness through the dualizing complex of the fibre. Same chapter, Lemma labelled lemma-gorenstein, cited at proof lines 6749-6751. Gives the equivalence of Gorensteinness and invertibility on the fibre.
\item Lifting invertibility from the point to a neighbourhood. More on Algebra, Lemma labelled more-algebra-lemma-lift-bounded-pseudo-coherent-to-perfect, cited at proof lines 6753-6757 and 6759-6763. Moves between the dimension one condition at the point and invertibility nearby.
\end{itemize}\medskip
\ifdefined\hypertarget\hypertarget{sourcepointer-1}{}\fi
\noindent\textit{Source reference.} Original reference 0F42: sources/stacks-duality/pinned-original-source.tex, source lines 3953--3953. Same-chapter pointer into the author environment situation at source line 3953; invoked at proof line None of the counted proof. It is retained only as a navigation pointer; every load-bearing step is inside the delivered span.\par\medskip
\ifdefined\hypertarget\hypertarget{sourcepointer-2}{}\fi
\noindent\textit{Source reference.} Original lemma 0BZZ: sources/stacks-duality/pinned-original-source.tex, source lines 5130--5130. Same-chapter pointer into the author environment lemma at source line 5130; invoked at proof line 6737 of the counted proof. It is retained only as a navigation pointer; every load-bearing step is inside the delivered span.\par\medskip
\ifdefined\hypertarget\hypertarget{sourcepointer-3}{}\fi
\noindent\textit{Source reference.} Original lemma 0BFQ: sources/stacks-duality/pinned-original-source.tex, source lines 6400--6400. Same-chapter pointer into the author environment lemma at source line 6400; invoked at proof line 6755 of the counted proof. It is retained only as a navigation pointer; every load-bearing step is inside the delivered span.\par\medskip
\ifdefined\hypertarget\hypertarget{sourcepointer-4}{}\fi
\noindent\textit{Source reference.} Original lemma 0BYP: sources/stacks-chow/tags.txt, source lines 15022--15022. Cross-chapter standard prerequisite located in the project tag list staged in the local snapshot; invoked at proof line 6758 of the counted proof. It is retained only as a navigation pointer; every load-bearing step is inside the delivered span.\par\medskip
\begin{lemma}
\label{lemma-affine-flat-Noetherian-gorenstein}
In Situation \ref{situation-shriek} let $f : X \to Y$ be a morphism
of $\textit{FTS}_S$. Let $x \in X$. If $f$ is flat, then
the following are equivalent
\begin{enumerate}
\item $f$ is Gorenstein at $x$,
\item $f^!\mathcal{O}_Y$ is isomorphic to an invertible object
in a neighbourhood of $x$.
\end{enumerate}
In particular, the set of points where $f$ is Gorenstein is
open in $X$.
\end{lemma}

\begin{proof}
Set $\omega^\bullet = f^!\mathcal{O}_Y$. By
Lemma \ref{lemma-relative-dualizing-fibres}
this is a bounded complex with coherent cohomology
sheaves whose derived restriction $Lh^*\omega^\bullet$
to the fibre $X_y$ is a dualizing complex on $X_y$.
Denote $i : x \to X_y$ the inclusion of a point.
Then the following are equivalent
\begin{enumerate}
\item $f$ is Gorenstein at $x$,
\item $\mathcal{O}_{X_y, x}$ is Gorenstein,
\item $Lh^*\omega^\bullet$ is invertible in a neighbourhood of $x$,
\item $Li^* Lh^* \omega^\bullet$ has exactly one nonzero
cohomology of dimension $1$ over $\kappa(x)$,
\item $L(h \circ i)^* \omega^\bullet$ has exactly one nonzero
cohomology of dimension $1$ over $\kappa(x)$,
\item $\omega^\bullet$ is invertible in a neighbourhood of $x$.
\end{enumerate}
The equivalence of (1) and (2) is by definition (as $f$ is flat).
The equivalence of (2) and (3) follows from
Lemma \ref{lemma-gorenstein}.
The equivalence of (3) and (4) follows from
More on Algebra, Lemma
\ref{more-algebra-lemma-lift-bounded-pseudo-coherent-to-perfect}.
The equivalence of (4) and (5) holds because
$Li^* Lh^* = L(h \circ i)^*$.
The equivalence of (5) and (6) holds by
More on Algebra, Lemma
\ref{more-algebra-lemma-lift-bounded-pseudo-coherent-to-perfect}.
Thus the lemma is clear.
\end{proof}

\end{document}
